\section{Related Work: Graph Anomaly Detection}
\label{sec:related-gad}

\paragraph{From node scores to learned structural detectors.} Graph anomaly detection (GAD) matured from classical outlier methods operating on graph-derived features~\cite{zhao2019pyod,liu2022pygod} into deep detectors with dedicated benchmarks: BOND standardized unsupervised outlier-node detection on attributed graphs~\cite{liu2022bond}, and GADBench did the same for supervised detection, showing simple supervised baselines are embarrassingly strong~\cite{tang2023gadbench}. A structural thread then rethought what signal detectors should model: B-GAD links anomalousness to a homophily gap and resists heterophily~\cite{tang2022rethinking}; GDN targets heterophily through graph spectrum alignment~\cite{gao2023addressing}; GAD-NR reconstructs neighborhoods rather than node features to expose anomalous connectivity~\cite{roy2024gadnr}; ADA-GAD denoises the anomaly signal itself inside the autoencoder~\cite{he2024adagad}; SmoothGNN smooths anomalies out of message passing to sharpen their contrast~\cite{dong2025smoothgnn}; and dynamic-spectral detectors adapt filters to evolving structure~\cite{zheng2025dynamic}. These advances share an assumption that does not hold on agent traces: a large, mostly homogeneous graph where anomalies are rare nodes. Agent trajectory graphs are small, fully typed, temporally ordered, and anomalous \emph{as wholes}, which is the regime of graph-level detection.

\paragraph{Graph-level anomaly detection.} Graph-level GAD scores whole graphs against a distribution of normal graphs, surveyed comprehensively by~\cite{qiao2024deep}. The newest methods exploit language-aligned semantics and foundation-model transfer: GLASS aligns graph structure with language priors under a spherical scoring rule for transferable detection~\cite{wang2026glass}; AnomalyGFM and OWLEYE target zero- and few-shot anomaly detection across domains~\cite{qiao2025anomalygfm,zheng2026owleye}; and AD-LLM benchmarks LLMs as anomaly detectors on tabular and graph tasks~\cite{yang2024adllm}. What is striking is the testbed monoculture: GLAD benchmarks are built from molecules, program graphs, and small citation clusters---graphs without execution semantics, without sensitive operations, and without an adversarial generative process producing the anomalies. Agent trajectories break all three assumptions: nodes are typed actions and artifacts, edges encode argument flow under an adversary, and the anomaly families are generated by replaying actual attacks. Our benchmark therefore gives GLAD a genuinely new testbed class, and conversely supplies the agent-safety community with detector families that exploit exactly the structure that step-level guardrails ignore.

\paragraph{Adversarial pressure on graph detectors.} GAD methods are also attacked, and the attack literature maps one-to-one onto our robustness protocol. Node injection attacks insert malicious nodes whose attributes camouflage them into the normal class~\cite{ju2022blackbox,tao2022adversarial}; multi-target variants coordinate injected subgraphs against fraud detectors~\cite{choi2024unveiling}; and the newest work uses LLMs to generate text-level attributes for injected nodes that evade LLM-enhanced detectors~\cite{lyu2025navigating,ma2025graphtextack}. This is precisely the threat our mimicry protocol instantiates for trajectories: an adversary that splices benign-looking structure around the malicious payload. Prior GAD robustness work, however, injects nodes into \emph{large} graphs to poison a detector; in our setting the adversary controls the entire trajectory-generation process, which is strictly stronger, and our protocol measures the resulting degradation explicitly rather than leaving it as future work.
